% Generated from lessons/0009-1-numerical-series-convergence.html; edit the HTML source, then rerun the converter.
\section{数项级数的收敛性}
\lessonmark{数项级数的收敛性}

这一节只解决一个基础问题：无限多个数相加到底是什么意思，什么时候有“和”。考试里要把级数看成部分和数列的极限，而不是把无穷多个数真的一次性加完。

\subsection*{本节主线}

\begin{conceptbox}
\textbf{定义级数}
级数 \(\sum x_n\) 的本质是部分和数列 \(S_n=x_1+\cdots+x_n\)。
\end{conceptbox}

\begin{conceptbox}
\textbf{判断收敛}
看 \(S_n\) 是否有有限极限。若有，极限就是级数的和。
\end{conceptbox}

\begin{conceptbox}
\textbf{计算套路}
几何级数、裂项相消、分组、参数化求和，是这一节的主要工具。
\end{conceptbox}

\begin{notebox}
一句话记忆：级数收敛不是“项越来越小”，而是“部分和越来越稳定”。
\end{notebox}

\subsection*{一、定义：级数是部分和的极限}

给定数列 \(x_1,x_2,\ldots,x_n,\ldots\)，把形式和

\[
      x_1+x_2+\cdots+x_n+\cdots
      =
      \sum_{n=1}^{\infty}x_n
    \]

称为一个数项级数。它的第 \(n\) 个部分和是

\[
      S_n=\sum_{k=1}^{n}x_k.
    \]

若部分和数列 \(\{S_n\}\) 收敛，且

\[
      \lim_{n\to\infty}S_n=S,
    \]

则称级数 \(\sum_{n=1}^{\infty}x_n\) 收敛，并称 \(S\) 为它的和，记作

\[
      \sum_{n=1}^{\infty}x_n=S.
    \]

若 \(\{S_n\}\) 不收敛，则级数发散。

\begin{notebox}
所以判断级数敛散，第一原则永远是：研究 \(S_n\)，而不是只盯着 \(x_n\)。
\end{notebox}

\subsection*{二、基本例子}

\subsubsection*{1. 等比级数}

若 \(|q|<1\)，则

\[
      \sum_{n=0}^{\infty}q^n=\frac{1}{1-q}.
    \]

证明只用有限等比求和：

\[
      S_n=1+q+\cdots+q^n=\frac{1-q^{n+1}}{1-q},
      \qquad q\ne1.
    \]

当 \(|q|<1\) 时 \(q^{n+1}\to0\)，于是 \(S_n\to \frac1{1-q}\)。若 \(|q|\ge1\)，通常由一般项不趋于 \(0\) 或部分和振荡推出发散。

\subsubsection*{2. \(1-1+1-1+\cdots\) 不收敛}

它的一部分部分和是 \(1\)，另一部分是 \(0\)，所以部分和数列没有极限。

\[
      S_1=1,\quad S_2=0,\quad S_3=1,\quad S_4=0,\ldots
    \]

\subsubsection*{3. \(p\)-级数是参照物}

\[
      \sum_{n=1}^{\infty}\frac1{n^p}
      \begin{cases}
      \text{收敛}, & p>1,\\
      \text{发散}, & 0<p\le1.
      \end{cases}
    \]

这一节先把它当作标准结论使用；后面正项级数会系统证明。考试中看到 \(\frac1{n^p}\)、\(\frac1{x^p}\)，要立刻想到它。

\subsection*{三、余项与必要条件}

\subsubsection*{余项}

若级数 \(\sum x_n\) 收敛，和为 \(S\)，则第 \(n\) 个余项是

\[
      r_n=\sum_{k=n+1}^{\infty}x_k=S-S_n.
    \]

因为 \(S_n\to S\)，所以

\[
      r_n\to0.
    \]

\subsubsection*{定理 9.1.1：收敛级数的一般项趋于零}

若 \(\sum_{n=1}^{\infty}x_n\) 收敛，则

\[
      \lim_{n\to\infty}x_n=0.
    \]

\textbf{证明思路：}设 \(S_n=x_1+\cdots+x_n\)，收敛时 \(S_n\to S\)，也有 \(S_{n-1}\to S\)。而

\[
      x_n=S_n-S_{n-1}\to S-S=0.
    \]

\begin{notebox}
这个条件只能用来快速判发散：若 \(x_n\not\to0\)，则 \(\sum x_n\) 一定发散。反过来不成立，例如 \(\sum\frac1n\) 的一般项趋于 \(0\)，但级数发散。
\end{notebox}

\subsection*{四、线性性质}

\subsubsection*{定理 9.1.2：收敛级数可线性组合}

若

\[
      \sum_{n=1}^{\infty}a_n=A,\qquad
      \sum_{n=1}^{\infty}b_n=B,
    \]

且 \(\alpha,\beta\) 为常数，则

\[
      \sum_{n=1}^{\infty}(\alpha a_n+\beta b_n)
      =
      \alpha A+\beta B.
    \]

\textbf{证明思路：}取部分和，利用有限和的线性性，再令 \(n\to\infty\)：

\[
      \sum_{k=1}^{n}(\alpha a_k+\beta b_k)
      =
      \alpha\sum_{k=1}^{n}a_k+\beta\sum_{k=1}^{n}b_k.
    \]

\begin{notebox}
注意前提：两个原级数都已经收敛。不能随便把两个发散级数拆开再相减。
\end{notebox}

\subsection*{五、加括号：只能保护已收敛的级数}

\subsubsection*{定理 9.1.3：收敛级数任意加括号，和不变}

若级数 \(\sum x_n\) 收敛，把它的项任意连续分组，例如

\[
      (x_1+\cdots+x_{n_1})+(x_{n_1+1}+\cdots+x_{n_2})+\cdots
    \]

所得的新级数仍收敛，且和与原级数相同。

\textbf{证明思路：}新级数的部分和只是原级数部分和数列 \(\{S_n\}\) 的一个子列。原数列收敛时，它的子列也收敛到同一个极限。

\begin{notebox}
反过来不成立。发散级数可能通过加括号得到看起来“收敛”的结果，例如 \(1-1+1-1+\cdots\) 可以被括成 \((1-1)+(1-1)+\cdots=0\)，也可以括成 \(1+(-1+1)+\cdots=1\)。这说明：不能靠加括号拯救一个本来发散的级数。
\end{notebox}

\subsection*{六、考试常用计算模板}

\begin{center}
\small
\begin{tabularx}{\linewidth}{|>{\raggedright\arraybackslash}p{0.18\linewidth}|>{\raggedright\arraybackslash}p{0.42\linewidth}|>{\raggedright\arraybackslash}X|}
\hline
\textbf{类型} & \textbf{动作} & \textbf{典型形状} \\ \hline
一般项不趋零 & 直接判发散 & \(\displaystyle \frac{2n}{3n+1}\to\frac23\) \\ \hline
等比级数 & 化成 \(\sum ar^n\) & \(\displaystyle \sum \frac{1}{2^n},\quad \sum e^{nx}\) \\ \hline
裂项相消 & 拆成相邻差 & \(\displaystyle \frac1{n(n+2)}=\frac12\left(\frac1n-\frac1{n+2}\right)\) \\ \hline
二阶差分 & 看作 \(a_{n+2}-2a_{n+1}+a_n\) & \(\sqrt{n+2}-2\sqrt{n+1}+\sqrt n\) \\ \hline
循环小数 & 整数部分加等比尾巴 & \((36.07360736\cdots)_8\) \\ \hline
\end{tabularx}
\end{center}

\subsection*{七、即时判断练习}

\begin{quickcheck}
若 \(\lim_{n\to\infty}x_n\ne0\)，则 \(\sum x_n\) 怎么样？

\tcblower
\textbf{答案：}
\end{quickcheck}

\begin{quickcheck}
若 \(\sum x_n\) 收敛，则任意连续加括号后得到的新级数？

\tcblower
\textbf{答案：}
\end{quickcheck}

\begin{quickcheck}
\(\sum_{n=1}^{\infty}\frac1{\sqrt[4]{n}}\) 的敛散性是？

\tcblower
\textbf{答案：}
\end{quickcheck}

\subsection*{八、课后题训练}

下面按教材第九章第 1 节习题整理。题面完整保留；提示只给入口，写作业时建议先独立算一遍，再展开提示核对方向。

\begin{exercisebox}
\textbf{习题 1 \badge{敛散性与求和}}

讨论下列级数的敛散性。收敛的话，试求出级数之和：

\[
      \begin{aligned}
      &(1)\ \sum_{n=1}^{\infty}\frac1{n(n+2)},\qquad
      &&(2)\ \sum_{n=1}^{\infty}\frac{2n}{3n+1},\\
      &(3)\ \sum_{n=1}^{\infty}\frac1{n(n+1)(n+2)},\qquad
      &&(4)\ \sum_{n=1}^{\infty}\left(\frac1{2^n}-\frac1{3^n}\right),\\
      &(5)\ \sum_{n=1}^{\infty}\frac1{\sqrt[4]{n}},\qquad
      &&(6)\ \sum_{n=1}^{\infty}\frac{5^{n-1}+4^{n+1}}{3^{2n}},\\
      &(7)\ \sum_{n=1}^{\infty}\left(\sqrt{n+2}-2\sqrt{n+1}+\sqrt n\right),\qquad
      &&(8)\ \sum_{n=1}^{\infty}\frac{2n-1}{3^n},\\
      &(9)\ \sum_{n=0}^{\infty}q^n\cos n\theta\quad (|q|<1).
      \end{aligned}
      \]

\begin{hintbox}
\begin{itemize}
 \item (1) 用 \(\displaystyle \frac1{n(n+2)}=\frac12\left(\frac1n-\frac1{n+2}\right)\)，和为 \(\frac34\)。
 \item (2) 一般项趋于 \(\frac23\)，所以发散。
 \item (3) 用 \(\displaystyle \frac1{n(n+1)(n+2)}=\frac12\left[\frac1{n(n+1)}-\frac1{(n+1)(n+2)}\right]\)，和为 \(\frac14\)。
 \item (4) 两个等比级数相减，和为 \(\frac12\)。
 \item (5) 是 \(p=\frac14\) 的 \(p\)-级数，发散。
 \item (6) 拆成两个等比级数，和为 \(\frac{69}{20}\)。
 \item (7) 看成 \(a_{n+2}-2a_{n+1}+a_n\)，其中 \(a_n=\sqrt n\)，和为 \(1-\sqrt2\)。
 \item (8) 用 \(\sum nr^n=\frac{r}{(1-r)^2}\)，和为 \(1\)。
 \item (9) 可取实部：\(\displaystyle \sum_{n=0}^{\infty}(qe^{i\theta})^n=\frac1{1-qe^{i\theta}}\)，结果为 \(\displaystyle \frac{1-q\cos\theta}{1-2q\cos\theta+q^2}\)。
 \end{itemize}
\end{hintbox}
\end{exercisebox}

\begin{exercisebox}
\textbf{习题 2 \badge{含参数等比级数}}

确定 \(x\) 的范围，使下列级数收敛：

\[
      (1)\ \sum_{n=1}^{\infty}\frac1{(1-x)^n},\qquad
      (2)\ \sum_{n=1}^{\infty}e^{nx},\qquad
      (3)\ \sum_{n=1}^{\infty}x^n(1-x).
      \]

\begin{hintbox}
全部先识别成等比级数。

\begin{itemize}
 \item (1) 公比为 \(\frac1{1-x}\)，需 \(\left|\frac1{1-x}\right|<1\)，即 \(x<0\) 或 \(x>2\)。
 \item (2) 公比为 \(e^x\)，需 \(e^x<1\)，即 \(x<0\)。
 \item (3) 当 \(|x|<1\) 收敛；当 \(x=1\) 时每项为 \(0\)，也收敛；其余发散。因此范围是 \(-1<x\le1\)。
 \end{itemize}
\end{hintbox}
\end{exercisebox}

\begin{exercisebox}
\textbf{习题 3 \badge{进位制循环小数}}

求八进制无限循环小数

\[
        (36.0736073607\cdots)_8
      \]

的值。

\begin{hintbox}
整数部分是 \((36)_8=30\)。小数部分的循环节是 \(0736\)，所以

\[
          0.(0736)_8=\frac{(0736)_8}{8^4-1}
          =
          \frac{478}{4095}.
        \]

因此总值为 \(\displaystyle 30+\frac{478}{4095}=\frac{123328}{4095}\)。
\end{hintbox}
\end{exercisebox}

\begin{exercisebox}
\textbf{习题 4 \badge{积分与级数}}

设

\[
        x_n=\int_0^1 x^2(1-x)^n\,dx,
      \]

求级数 \(\displaystyle \sum_{n=1}^{\infty}x_n\) 的和。

\begin{hintbox}
先对有限部分和求和：

\[
        \sum_{n=1}^{N}x_n
        =
        \int_0^1 x^2\sum_{n=1}^{N}(1-x)^n\,dx.
        \]

令 \(N\to\infty\)，几何和给出 \(\sum_{n=1}^{\infty}(1-x)^n=\frac{1-x}{x}\)，所以

\[
        \sum_{n=1}^{\infty}x_n
        =
        \int_0^1 x(1-x)\,dx
        =
        \frac16.
        \]
\end{hintbox}
\end{exercisebox}

\begin{exercisebox}
\textbf{习题 5 \badge{几何面积与级数}}

设抛物线

\[
        l_n:\ y=nx^2+\frac1n,\qquad
        l_n':\ y=(n+1)x^2+\frac1{n+1}
      \]

的交点的横坐标的绝对值为 \(a_n\;(n=1,2,\ldots)\)。

\begin{enumerate}
 \item 求抛物线 \(l_n\) 与 \(l_n'\) 所围成的平面图形的面积 \(S_n\)；
 \item 求级数 \(\displaystyle \sum_{n=1}^{\infty}\frac{S_n}{a_n}\) 的和。
 \end{enumerate}

\begin{hintbox}
交点满足

\[
        nx^2+\frac1n=(n+1)x^2+\frac1{n+1},
        \qquad
        a_n^2=\frac1{n(n+1)}.
        \]

在 \([-a_n,a_n]\) 上，\(l_n\) 在上，面积为

\[
        S_n=\int_{-a_n}^{a_n}(a_n^2-x^2)\,dx
        =
        \frac43a_n^3.
        \]

于是

\[
        \frac{S_n}{a_n}=\frac43a_n^2=\frac4{3n(n+1)},
        \qquad
        \sum_{n=1}^{\infty}\frac{S_n}{a_n}
        =
        \frac43.
        \]
\end{hintbox}
\end{exercisebox}

\subsection*{九、这一节的选题建议}

\begin{conceptbox}
\textbf{必做}
习题 1 的 (1)(2)(3)(7)(9)，习题 2，习题 4。
\end{conceptbox}

\begin{conceptbox}
\textbf{好题}
习题 5 把定积分面积和级数合在一起，适合放进好题集。
\end{conceptbox}

\begin{conceptbox}
\textbf{检查点}
每题先问：一般项是否趋零？能否写部分和？能否裂项或化等比？
\end{conceptbox}

来源参考：教材《数学分析》第三版下册第九章 §1（数项级数的收敛性）。如果哪一道题卡住，可以直接在页面中选中对应题目问我。
